% preamble.tex -- pengaturan bersama seluruh gambar TikZ.
% Dipanggil dengan % preamble.tex -- pengaturan bersama seluruh gambar TikZ.
% Dipanggil dengan % preamble.tex -- pengaturan bersama seluruh gambar TikZ.
% Dipanggil dengan % preamble.tex -- pengaturan bersama seluruh gambar TikZ.
% Dipanggil dengan \input{preamble} tepat setelah \documentclass{standalone}.
%
% Font disamakan dengan template.cls (newtxtext/newtxmath) supaya teks di
% dalam gambar menyatu dengan teks naskah. Tanpa ini, angka pada sumbu grafik
% terlihat berbeda dari angka pada paragraf di sebelahnya.

\usepackage[T1]{fontenc}
\usepackage{newtxtext,newtxmath}
\usepackage{tikz}
\usepackage{pgfplots}
\pgfplotsset{compat=1.18}
\usetikzlibrary{calc,positioning,arrows.meta,decorations.pathreplacing,fit,backgrounds,shapes.geometric}

% Warna konsisten antar gambar.
\definecolor{aksen}{RGB}{0,90,160}      % biru: kurva atau objek utama kedua
\definecolor{aksendua}{RGB}{170,60,20}  % jingga gelap: objek ketiga
\definecolor{netral}{gray}{0.45}

% Gaya baku untuk diagram alur.
\tikzset{
  tahap/.style={
    rectangle, rounded corners=2pt, draw=black!55, fill=gray!8,
    align=center, font=\footnotesize, inner sep=5pt,
  },
  tahapaksen/.style={
    tahap, draw=aksen, fill=aksen!8,
  },
  panah/.style={
    -{Latex[length=5pt,width=4pt]}, line width=0.7pt, draw=black!60,
  },
  panahaksen/.style={
    panah, draw=aksen,
  },
  kelompok/.style={
    rectangle, rounded corners=3pt, draw=netral, dashed, line width=0.5pt,
    inner sep=7pt,
  },
}

% Gaya sumbu baku: garis tengah, grid tipis, label kecil.
\pgfplotsset{
  bakuaxis/.style={
    axis lines=left,
    axis line style={gray!70},
    tick label style={font=\footnotesize},
    label style={font=\small},
    legend style={font=\footnotesize, draw=gray!50, fill=white,
                  fill opacity=0.9, text opacity=1, row sep=1pt},
    grid=both,
    grid style={line width=0.2pt, draw=gray!18},
    every axis plot/.append style={line width=1.0pt},
  },
  bakubar/.style={
    ybar,
    bar width=\barw,
    fill=gray!35,
    draw=black!70,
    line width=0.5pt,
  },
}

\newcommand{\barw}{7pt}

% ------------------------------------------------------------------
% Grid matriks. Dipakai untuk gambar konvolusi, pooling, padding, dan
% perhitungan manual lain yang menampilkan nilai per sel.
%
%   \gridmat{<ukuran sel>}{<gaya tambahan sel>}{{{1,2},{3,4}}}
%
% Sel pertama menempati kuadran kanan-bawah titik (0,0), sehingga baris
% bertambah ke bawah dan kolom bertambah ke kanan seperti indeks matriks.
% Catatan: pgfmath bekerja dalam satuan pt dengan batas ~16384, jadi hindari
% hasil antara yang besar pada nilai sel.
% ------------------------------------------------------------------
\tikzset{
  sel/.style={draw=black!55, fill=white, line width=0.35pt},
  selaksen/.style={sel, fill=aksen!18, draw=aksen},
  seldua/.style={sel, fill=aksendua!18, draw=aksendua},
  selnol/.style={sel, fill=gray!12, draw=black!35},
}
% Ukuran huruf isi sel. Di-\renewcommand pada gambar yang selnya rapat.
\newcommand{\gridfont}{\footnotesize}
\newcommand{\gridmat}[3]{%
  \foreach \gridrow [count=\gr from 0] in #3 {%
    \foreach \gridval [count=\gc from 0] in \gridrow {%
      \draw[sel,#2] (\gc*#1, -\gr*#1) rectangle ++(#1,-#1);
      \node[font=\gridfont] at (\gc*#1+0.5*#1, -\gr*#1-0.5*#1) {\gridval};
    }%
  }%
}

% Sorotan window pada grid: #1 ukuran sel, #2 kolom awal (0-basis),
% #3 baris awal (0-basis), #4 lebar window, #5 gaya garis.
%   \window{0.6}{1}{1}{3}{draw=aksen, line width=1.2pt}
\newcommand{\window}[5]{%
  \draw[#5] (#2*#1, -#3*#1) rectangle ++(#4*#1, -#4*#1);}

% Label panel (a)/(b)/(c) di bawah sebuah axis bernama.
% Anchor "below south" milik pgfplots berada di bawah seluruh elemen axis,
% termasuk label sumbu, sehingga label panel tidak pernah bertabrakan dengannya.
\newcommand{\panellabel}[2]{%
  \node[font=\small, anchor=north, yshift=-1.5mm] at (#1.below south) {#2};}

% Keterangan bersama di bawah dua axis bernama.
\newcommand{\keteranganbersama}[3]{%
  \node[font=\scriptsize, align=center, anchor=north, yshift=-10mm]
    at ($(#1.below south)!0.5!(#2.below south)$) {#3};}
 tepat setelah \documentclass{standalone}.
%
% Font disamakan dengan template.cls (newtxtext/newtxmath) supaya teks di
% dalam gambar menyatu dengan teks naskah. Tanpa ini, angka pada sumbu grafik
% terlihat berbeda dari angka pada paragraf di sebelahnya.

\usepackage[T1]{fontenc}
\usepackage{newtxtext,newtxmath}
\usepackage{tikz}
\usepackage{pgfplots}
\pgfplotsset{compat=1.18}
\usetikzlibrary{calc,positioning,arrows.meta,decorations.pathreplacing,fit,backgrounds,shapes.geometric}

% Warna konsisten antar gambar.
\definecolor{aksen}{RGB}{0,90,160}      % biru: kurva atau objek utama kedua
\definecolor{aksendua}{RGB}{170,60,20}  % jingga gelap: objek ketiga
\definecolor{netral}{gray}{0.45}

% Gaya baku untuk diagram alur.
\tikzset{
  tahap/.style={
    rectangle, rounded corners=2pt, draw=black!55, fill=gray!8,
    align=center, font=\footnotesize, inner sep=5pt,
  },
  tahapaksen/.style={
    tahap, draw=aksen, fill=aksen!8,
  },
  panah/.style={
    -{Latex[length=5pt,width=4pt]}, line width=0.7pt, draw=black!60,
  },
  panahaksen/.style={
    panah, draw=aksen,
  },
  kelompok/.style={
    rectangle, rounded corners=3pt, draw=netral, dashed, line width=0.5pt,
    inner sep=7pt,
  },
}

% Gaya sumbu baku: garis tengah, grid tipis, label kecil.
\pgfplotsset{
  bakuaxis/.style={
    axis lines=left,
    axis line style={gray!70},
    tick label style={font=\footnotesize},
    label style={font=\small},
    legend style={font=\footnotesize, draw=gray!50, fill=white,
                  fill opacity=0.9, text opacity=1, row sep=1pt},
    grid=both,
    grid style={line width=0.2pt, draw=gray!18},
    every axis plot/.append style={line width=1.0pt},
  },
  bakubar/.style={
    ybar,
    bar width=\barw,
    fill=gray!35,
    draw=black!70,
    line width=0.5pt,
  },
}

\newcommand{\barw}{7pt}

% ------------------------------------------------------------------
% Grid matriks. Dipakai untuk gambar konvolusi, pooling, padding, dan
% perhitungan manual lain yang menampilkan nilai per sel.
%
%   \gridmat{<ukuran sel>}{<gaya tambahan sel>}{{{1,2},{3,4}}}
%
% Sel pertama menempati kuadran kanan-bawah titik (0,0), sehingga baris
% bertambah ke bawah dan kolom bertambah ke kanan seperti indeks matriks.
% Catatan: pgfmath bekerja dalam satuan pt dengan batas ~16384, jadi hindari
% hasil antara yang besar pada nilai sel.
% ------------------------------------------------------------------
\tikzset{
  sel/.style={draw=black!55, fill=white, line width=0.35pt},
  selaksen/.style={sel, fill=aksen!18, draw=aksen},
  seldua/.style={sel, fill=aksendua!18, draw=aksendua},
  selnol/.style={sel, fill=gray!12, draw=black!35},
}
% Ukuran huruf isi sel. Di-\renewcommand pada gambar yang selnya rapat.
\newcommand{\gridfont}{\footnotesize}
\newcommand{\gridmat}[3]{%
  \foreach \gridrow [count=\gr from 0] in #3 {%
    \foreach \gridval [count=\gc from 0] in \gridrow {%
      \draw[sel,#2] (\gc*#1, -\gr*#1) rectangle ++(#1,-#1);
      \node[font=\gridfont] at (\gc*#1+0.5*#1, -\gr*#1-0.5*#1) {\gridval};
    }%
  }%
}

% Sorotan window pada grid: #1 ukuran sel, #2 kolom awal (0-basis),
% #3 baris awal (0-basis), #4 lebar window, #5 gaya garis.
%   \window{0.6}{1}{1}{3}{draw=aksen, line width=1.2pt}
\newcommand{\window}[5]{%
  \draw[#5] (#2*#1, -#3*#1) rectangle ++(#4*#1, -#4*#1);}

% Label panel (a)/(b)/(c) di bawah sebuah axis bernama.
% Anchor "below south" milik pgfplots berada di bawah seluruh elemen axis,
% termasuk label sumbu, sehingga label panel tidak pernah bertabrakan dengannya.
\newcommand{\panellabel}[2]{%
  \node[font=\small, anchor=north, yshift=-1.5mm] at (#1.below south) {#2};}

% Keterangan bersama di bawah dua axis bernama.
\newcommand{\keteranganbersama}[3]{%
  \node[font=\scriptsize, align=center, anchor=north, yshift=-10mm]
    at ($(#1.below south)!0.5!(#2.below south)$) {#3};}
 tepat setelah \documentclass{standalone}.
%
% Font disamakan dengan template.cls (newtxtext/newtxmath) supaya teks di
% dalam gambar menyatu dengan teks naskah. Tanpa ini, angka pada sumbu grafik
% terlihat berbeda dari angka pada paragraf di sebelahnya.

\usepackage[T1]{fontenc}
\usepackage{newtxtext,newtxmath}
\usepackage{tikz}
\usepackage{pgfplots}
\pgfplotsset{compat=1.18}
\usetikzlibrary{calc,positioning,arrows.meta,decorations.pathreplacing,fit,backgrounds,shapes.geometric}

% Warna konsisten antar gambar.
\definecolor{aksen}{RGB}{0,90,160}      % biru: kurva atau objek utama kedua
\definecolor{aksendua}{RGB}{170,60,20}  % jingga gelap: objek ketiga
\definecolor{netral}{gray}{0.45}

% Gaya baku untuk diagram alur.
\tikzset{
  tahap/.style={
    rectangle, rounded corners=2pt, draw=black!55, fill=gray!8,
    align=center, font=\footnotesize, inner sep=5pt,
  },
  tahapaksen/.style={
    tahap, draw=aksen, fill=aksen!8,
  },
  panah/.style={
    -{Latex[length=5pt,width=4pt]}, line width=0.7pt, draw=black!60,
  },
  panahaksen/.style={
    panah, draw=aksen,
  },
  kelompok/.style={
    rectangle, rounded corners=3pt, draw=netral, dashed, line width=0.5pt,
    inner sep=7pt,
  },
}

% Gaya sumbu baku: garis tengah, grid tipis, label kecil.
\pgfplotsset{
  bakuaxis/.style={
    axis lines=left,
    axis line style={gray!70},
    tick label style={font=\footnotesize},
    label style={font=\small},
    legend style={font=\footnotesize, draw=gray!50, fill=white,
                  fill opacity=0.9, text opacity=1, row sep=1pt},
    grid=both,
    grid style={line width=0.2pt, draw=gray!18},
    every axis plot/.append style={line width=1.0pt},
  },
  bakubar/.style={
    ybar,
    bar width=\barw,
    fill=gray!35,
    draw=black!70,
    line width=0.5pt,
  },
}

\newcommand{\barw}{7pt}

% ------------------------------------------------------------------
% Grid matriks. Dipakai untuk gambar konvolusi, pooling, padding, dan
% perhitungan manual lain yang menampilkan nilai per sel.
%
%   \gridmat{<ukuran sel>}{<gaya tambahan sel>}{{{1,2},{3,4}}}
%
% Sel pertama menempati kuadran kanan-bawah titik (0,0), sehingga baris
% bertambah ke bawah dan kolom bertambah ke kanan seperti indeks matriks.
% Catatan: pgfmath bekerja dalam satuan pt dengan batas ~16384, jadi hindari
% hasil antara yang besar pada nilai sel.
% ------------------------------------------------------------------
\tikzset{
  sel/.style={draw=black!55, fill=white, line width=0.35pt},
  selaksen/.style={sel, fill=aksen!18, draw=aksen},
  seldua/.style={sel, fill=aksendua!18, draw=aksendua},
  selnol/.style={sel, fill=gray!12, draw=black!35},
}
% Ukuran huruf isi sel. Di-\renewcommand pada gambar yang selnya rapat.
\newcommand{\gridfont}{\footnotesize}
\newcommand{\gridmat}[3]{%
  \foreach \gridrow [count=\gr from 0] in #3 {%
    \foreach \gridval [count=\gc from 0] in \gridrow {%
      \draw[sel,#2] (\gc*#1, -\gr*#1) rectangle ++(#1,-#1);
      \node[font=\gridfont] at (\gc*#1+0.5*#1, -\gr*#1-0.5*#1) {\gridval};
    }%
  }%
}

% Sorotan window pada grid: #1 ukuran sel, #2 kolom awal (0-basis),
% #3 baris awal (0-basis), #4 lebar window, #5 gaya garis.
%   \window{0.6}{1}{1}{3}{draw=aksen, line width=1.2pt}
\newcommand{\window}[5]{%
  \draw[#5] (#2*#1, -#3*#1) rectangle ++(#4*#1, -#4*#1);}

% Label panel (a)/(b)/(c) di bawah sebuah axis bernama.
% Anchor "below south" milik pgfplots berada di bawah seluruh elemen axis,
% termasuk label sumbu, sehingga label panel tidak pernah bertabrakan dengannya.
\newcommand{\panellabel}[2]{%
  \node[font=\small, anchor=north, yshift=-1.5mm] at (#1.below south) {#2};}

% Keterangan bersama di bawah dua axis bernama.
\newcommand{\keteranganbersama}[3]{%
  \node[font=\scriptsize, align=center, anchor=north, yshift=-10mm]
    at ($(#1.below south)!0.5!(#2.below south)$) {#3};}
 tepat setelah \documentclass{standalone}.
%
% Font disamakan dengan template.cls (newtxtext/newtxmath) supaya teks di
% dalam gambar menyatu dengan teks naskah. Tanpa ini, angka pada sumbu grafik
% terlihat berbeda dari angka pada paragraf di sebelahnya.

\usepackage[T1]{fontenc}
\usepackage{newtxtext,newtxmath}
\usepackage{tikz}
\usepackage{pgfplots}
\pgfplotsset{compat=1.18}
\usetikzlibrary{calc,positioning,arrows.meta,decorations.pathreplacing,fit,backgrounds,shapes.geometric}

% Warna konsisten antar gambar.
\definecolor{aksen}{RGB}{0,90,160}      % biru: kurva atau objek utama kedua
\definecolor{aksendua}{RGB}{170,60,20}  % jingga gelap: objek ketiga
\definecolor{netral}{gray}{0.45}

% Gaya baku untuk diagram alur.
\tikzset{
  tahap/.style={
    rectangle, rounded corners=2pt, draw=black!55, fill=gray!8,
    align=center, font=\footnotesize, inner sep=5pt,
  },
  tahapaksen/.style={
    tahap, draw=aksen, fill=aksen!8,
  },
  panah/.style={
    -{Latex[length=5pt,width=4pt]}, line width=0.7pt, draw=black!60,
  },
  panahaksen/.style={
    panah, draw=aksen,
  },
  kelompok/.style={
    rectangle, rounded corners=3pt, draw=netral, dashed, line width=0.5pt,
    inner sep=7pt,
  },
}

% Gaya sumbu baku: garis tengah, grid tipis, label kecil.
\pgfplotsset{
  bakuaxis/.style={
    axis lines=left,
    axis line style={gray!70},
    tick label style={font=\footnotesize},
    label style={font=\small},
    legend style={font=\footnotesize, draw=gray!50, fill=white,
                  fill opacity=0.9, text opacity=1, row sep=1pt},
    grid=both,
    grid style={line width=0.2pt, draw=gray!18},
    every axis plot/.append style={line width=1.0pt},
  },
  bakubar/.style={
    ybar,
    bar width=\barw,
    fill=gray!35,
    draw=black!70,
    line width=0.5pt,
  },
}

\newcommand{\barw}{7pt}

% ------------------------------------------------------------------
% Grid matriks. Dipakai untuk gambar konvolusi, pooling, padding, dan
% perhitungan manual lain yang menampilkan nilai per sel.
%
%   \gridmat{<ukuran sel>}{<gaya tambahan sel>}{{{1,2},{3,4}}}
%
% Sel pertama menempati kuadran kanan-bawah titik (0,0), sehingga baris
% bertambah ke bawah dan kolom bertambah ke kanan seperti indeks matriks.
% Catatan: pgfmath bekerja dalam satuan pt dengan batas ~16384, jadi hindari
% hasil antara yang besar pada nilai sel.
% ------------------------------------------------------------------
\tikzset{
  sel/.style={draw=black!55, fill=white, line width=0.35pt},
  selaksen/.style={sel, fill=aksen!18, draw=aksen},
  seldua/.style={sel, fill=aksendua!18, draw=aksendua},
  selnol/.style={sel, fill=gray!12, draw=black!35},
}
% Ukuran huruf isi sel. Di-\renewcommand pada gambar yang selnya rapat.
\newcommand{\gridfont}{\footnotesize}
\newcommand{\gridmat}[3]{%
  \foreach \gridrow [count=\gr from 0] in #3 {%
    \foreach \gridval [count=\gc from 0] in \gridrow {%
      \draw[sel,#2] (\gc*#1, -\gr*#1) rectangle ++(#1,-#1);
      \node[font=\gridfont] at (\gc*#1+0.5*#1, -\gr*#1-0.5*#1) {\gridval};
    }%
  }%
}

% Sorotan window pada grid: #1 ukuran sel, #2 kolom awal (0-basis),
% #3 baris awal (0-basis), #4 lebar window, #5 gaya garis.
%   \window{0.6}{1}{1}{3}{draw=aksen, line width=1.2pt}
\newcommand{\window}[5]{%
  \draw[#5] (#2*#1, -#3*#1) rectangle ++(#4*#1, -#4*#1);}

% Label panel (a)/(b)/(c) di bawah sebuah axis bernama.
% Anchor "below south" milik pgfplots berada di bawah seluruh elemen axis,
% termasuk label sumbu, sehingga label panel tidak pernah bertabrakan dengannya.
\newcommand{\panellabel}[2]{%
  \node[font=\small, anchor=north, yshift=-1.5mm] at (#1.below south) {#2};}

% Keterangan bersama di bawah dua axis bernama.
\newcommand{\keteranganbersama}[3]{%
  \node[font=\scriptsize, align=center, anchor=north, yshift=-10mm]
    at ($(#1.below south)!0.5!(#2.below south)$) {#3};}
